
\def\PUDcodigo{ECA211}

\if\COMPILAR0
   \def\PUDcodacad{04507.15}
\else
   \def\PUDcodacad{04506.16}
\fi

\def\PUDdisciplina{Eletrônica I}

\def\PUDsemestre{S3}

\def\PUDhoras{80}

%NOVOS ---------------------------------------------------
\def\PUDhorasTeoria{40}
\def\PUDhorasPratica{40}

\def\PUDinterECA{---}

\def\PUDatuaisECA{---}

\def\PUDrecursos{Projetor multimídia; Quadro branco e pincel; Aulas em laboratório de eletrônica.}

%----------------------------------------------------------

\def\PUDcreditos{4}

\def\PUDprerequisito{---}

\def\PUDpreacad{---}

\def\PUDobjetivos{Compreender o princípio de funcionamento dos dispositivos semicondutores básicos;  Compreender o funcionamento das portas lógicas; Identificar suas funções em circuitos lógicos combinacionais para solução de problemas lógicos; Descrever o funcionamento dos elementos de memória (flip-flop’s); Conhecer circuitos seqüenciais e conversões A/D e D/A.}

\def\PUDementa{Diodo semicondutor; Terminologia e Conceitos da Eletrônica Digital; Sistemas de Numeração Hexadecimal e Binário; Funções Lógicas e Álgebra Booleana; Circuitos Lógicos Combinacionais; Simplificação de Circuitos através da Álgebra de Boole; Simplificação de Circuitos através de  Mapas de Karnaugh;  Circuitos Somadores; Circuitos Codificadores e Decodificadores; Multiplexadores e Demultiplexadores; Flip-Flop; Registradores; Contadores; Circuitos Sequenciais; Conversão Analógico x Digital e Digital x Analógico.}

\def\PUDconteudo{ & Diodo Semicondutor: Estrutura atômica, processo de dopagem, junção PN, polarização do diodo, curva caracteristica e reta de carga do diodo, diodo emissor de luz - LED;  
\\ 
 & Funções Lógicas: Conversões de sistemas de numeração, circuitos lógicos combinacionais empregando portas lógicas básicas, tabela verdade e equivalência entre blocos lógicos; 
\\ 
 &  Projeto e Análise de Circuitos Lógicos: Circuitos lógicos combinacionais simples, teoremas, leis booleanas, projetar circuitos lógicos combinacionais a partir de situações diversas, simplificação utilizando a álgebra Booleana, simplificação utilizando mapas de Karnaugh. 
\\ 
 &  Circuitos Aritméticos: Circuitos aritméticos básicos, cálculos básicos, implementar circuitos lógicos aritméticos; 
\\ 
 &  Circuitos de Processamento de dados: Multiplexadores e demultiplexadores, circuitos codificadores e decodificadores e portas OUX; 
\\ 
 &  Principais Elementos de Memória: Flip-flop RS, flip-flop JK, flip-flop D e T, diagramas de tempo, registradores de deslocamento. 
\\ 
 &  Circuitos Seqüenciais: descrição de circuitos sequenciais, diagramas de transição de estado, contadores síncronos e assíncronos; 
\\ 
 &  Circuitos Conversores A/D e D/A: Circuitos conversores D/A, circuitos conversores A/D. }

\def\PUDmetodo{Aulas expositórias que podem ser teóricas e/ou práticas, onde as práticas no laboratório serão marcadas no decorrer da disciplina.}

\def\PUDavalia{A avaliação será feita com aplicação de provas teóricas e/ou práticas, além da possibilidade de inclusão de trabalhos e seminários no decorrer da disciplina.}

\def\PUDbibbasica{ & \href{www.biblioteca.ifce.edu.br/index.asp?codigo_sophia=24478}{Cruz}, Eduardo Cesar Alves.  Eletrônica Aplicada.  2ª ed.  São Paulo, SP: Érica, 2012.  295p.  ISBN: 9788536501505. 
\\ 
 & \href{www.biblioteca.ifce.edu.br/index.asp?codigo_sophia=24081}{Idoeta}, Ivan V.;  Francisco G. Capuano.  Elementos de Eletrônica Digital.  39º ed.  São Paulo, SP: Érica, 2007.  524p.  ISBN: 9788571940192. 
\\ 
 & \href{www.biblioteca.ifce.edu.br/index.asp?codigo_sophia=62882}{Tocci}, R. J.; Widner, N. S.; Moss, G. L.  Sistemas Digitais – Princípios e Aplicações.  11ª ed.  São Paulo, SP: Pearson Prentice Hall, 2011.  817p.  ISBN: 9788576059226. }

\def\PUDbibcomplementar{ & \href{www.biblioteca.ifce.edu.br/index.asp?codigo_sophia=24597}{Garcia}, P. A.; Martini, J. S. C.  Eletrônica Digital Teoria e Laboratório.  2ª ed.  São Paulo, SP: Érica, 2008.  184p.  ISBN: 9788536501093.
\\
 & \href{www.biblioteca.ifce.edu.br/index.asp?codigo_sophia=23875}{Malvino}, A. P.  Eletrônica Digital - Volume 1.  4ª ed.  São Paulo, SP: Pearson Makron Books, 1997.  670p.  ISBN: 9788534603782.
%\\ 
 %& MALVINO, A.  P.  Eletrônica Digital - Volume 2, 4ª Edição, Pearson Makron Books, São Paulo, 1997. 
\\ 
 & \href{www.biblioteca.ifce.edu.br/index.asp?codigo_sophia=24345}{Sedra}, A. S.  Microeletrônica. 5ª ed.  São Paulo, SP: Pearson Prentice Hall, 2007.  848p.  ISBN: 9788576050223.
\\
 & \href{www.biblioteca.ifce.edu.br/index.asp?codigo_sophia=55600}{Mariotto}, Paulo Antônio.  Análise de Circuitos Elétricos. E-book.  Pearson.  390p.  ISBN: 9788587918062.
\\
 & \href{www.biblioteca.ifce.edu.br/index.asp?codigo_sophia=59323}{Nilsson}, James William; Riedel, Susan A.  Circuitos elétricos.  E-book.  Pearson.  ISBN: 9788543004785. }
%\\ 
 %& BOYLESTAD, R. L.  Dispositivos eletrônicos e teoria de circuitos.  8ª Edição, Pearson Prentice Hall, São Paulo, 2004. 
%\\ 
 %&  BIGNELL, J.  W, DONOVAN R. , Eletrônica Digital, 5ª Edição, Cengage Learning.  2010.   ISBN 9788522107452
%\\ 
 %&  GARUE, S. , Eletrônica Digital - Circuitos e Tecnologia, 1ª Edição, Editora Hemus.  ISBN 9788528901405. }

\def\PUDautor{Fabrício Bandeira da Silva}

\def\PUDdata{2013-06-20}

\def\PUDrevisao{3}

\def\PUDrevisor{Fabrício Bandeira da Silva}

\def\PUDdatarevisao{2019-05-15}

\PUDcorpo
